\begin{table}[H]
    \centering
    \caption{Identificaci\'on y priorizaci\'on de Casos de Uso del Sistema (CUS)}
    \label{tab:cus_prioridades}
    \begin{tabularx}{\textwidth}{|>{\bfseries\raggedright\arraybackslash}p{2.0cm}|>{\raggedright\arraybackslash}p{3.2cm}|X|>{\centering\arraybackslash}p{2.2cm}|}
        \hline
        Código CUS & Nombre del Caso de Uso & Descripción Breve & Prioridad \\
        \hline
        CUS-01 & Registrar Venta & Permite al vendedor seleccionar productos, verificar stock en tiempo real y reservar existencias provisionalmente. & Alta \\
        \hline
        CUS-02 & Procesar Pago y Comprobante & Permite al cajero liquidar la venta en diversos métodos de pago y emitir la boleta o factura legal. & Alta \\
        \hline
        CUS-03 & Gestionar Abastecimiento & Permite al administrador evaluar stock crítico, cotizar con proveedores y formalizar órdenes de compra. & Media-Alta \\
        \hline
        CUS-04 & Registrar Ingreso de Mercancía & Permite al encargado de almacén validar bultos recibidos de órdenes de compra e incrementar el inventario. & Media-Alta \\
        \hline
        CUS-05 & Realizar Control de Inventario & Permite ejecutar auditorías físicas, registrar pérdidas/mermas y recalibrar el kardex del sistema. & Media \\
        \hline
        CUS-06 & Gestionar Devoluciones y Cambios & Permite procesar reclamos y garantías de clientes emitiendo notas de crédito o entregando reemplazos. & Media \\
        \hline
    \end{tabularx}
    \par\vspace{1ex}
    {\raggedright \small \emph{Nota.} Elaboración propia (2026). Clasificación de prioridades y alcance operativo de los CUS.\par}
\end{table}
